\documentclass[11pt,a4paper,oneside]{book}
\usepackage{mathscommun}
\setniveau{Seconde}
\renewcommand{\ReferenceProgramme}{Programme français de seconde GT -- rentrée 2026}
\renewcommand{\EditionDocument}{Édition du 18 octobre 2026}
\hypersetup{pdftitle={Mathématiques de seconde - Exercices et problèmes}}
\begin{document}
\titrecollection{Exercices et problèmes}{Chercher, représenter, calculer et démontrer.}{Cent quarante-quatre exercices, six par chapitre, et six problèmes de synthèse. Les approfondissements sont signalés. La numérotation correspond à celle du cours et de l'autre cahier.}
\frontmatter
\chapter*{Mode d'emploi}
Chercher sans consulter immédiatement le corrigé. Écrire les hypothèses, le domaine et les unités. Une réponse numérique doit être accompagnée de la méthode ; une preuve doit traiter toutes les valeurs autorisées. Sauf demande contraire, conserver les valeurs exactes et arrondir seulement à la fin.

Les questions suivent les notions du cours. Les six exercices du chapitre 24 et les autres questions marquées « facultatif » prolongent le programme. Les automatismes courts se trouvent dans le guide et se pratiquent régulièrement.
\tableofcontents
\mainmatter
\chapter{Langage, ensembles et preuves}
\section*{Exercice 1.1 -- Ensembles}
Dans $E=\{1,\ldots,10\}$, soit $A$ l'ensemble des multiples de $2$ et $B$ celui des multiples de $3$. Donner $A\cap B$, $A\cup B$, $\overline A$ et leurs cardinaux.

\section*{Exercice 1.2 -- Couples}
Pour $A=\{0,1\}$ et $B=\{a,b,c\}$, écrire $A\times B$. Les couples $(0,a)$ et $(a,0)$ désignent-ils la même chose ?

\section*{Exercice 1.3 -- Négations}
Nier : « $x>2$ », « $x\geq0$ et $x<5$ », « $x=-1$ ou $x=1$ ».

\section*{Exercice 1.4 -- Réciproque}
Prouver : si un entier est multiple de $6$, il est multiple de $3$. Écrire la réciproque et décider si elle est vraie.

\section*{Exercice 1.5 -- Une égalité d'ensembles}
Démontrer $A\cap(B\cup C)=(A\cap B)\cup(A\cap C)$ en raisonnant sur l'appartenance d'un élément.

\section*{Exercice 1.6 -- L'ordre des mots}
Comparer : « pour tout réel $x$, il existe un réel $y$ tel que $x+y=0$ » et « il existe un réel $y$ tel que, pour tout réel $x$, $x+y=0$ ».

\chapter{Entiers, divisibilité et fractions}
\section*{Exercice 2.1 -- Diviseurs}
Lister les diviseurs positifs de $36$ et déterminer si $84$ est multiple de $6$, de $8$ et de $7$.

\section*{Exercice 2.2 -- Dernier multiple}
Trouver le plus grand multiple de $13$ inférieur ou égal à $250$. Justifier qu'il n'en existe pas de plus grand sous cette borne.

\section*{Exercice 2.3 -- Fraction}
Réduire $630/462$ en une fraction irréductible.

\section*{Exercice 2.4 -- Parité}
Prouver que la somme de deux entiers impairs est paire et que leur produit est impair.

\section*{Exercice 2.5 -- Lots}
On dispose de $96$ billes rouges et $144$ bleues. Former le plus grand nombre de lots identiques sans reste.

\section*{Exercice 2.6 -- Un reste invariant}
Approfondissement : montrer que, pour tout entier impair $n$, $n^2-1$ est multiple de $8$.

\chapter{Réels, intervalles et approximations}
\section*{Exercice 3.1 -- Classer}
Donner le plus petit ensemble parmi $\N,\Z,\D,\Q,\R$ contenant chacun de $-4$, $7/8$, $2/3$, $\sqrt{25}$ et $\sqrt3$. On admet ici l'irrationalité de $\sqrt3$.

\section*{Exercice 3.2 -- Intervalles}
Pour $I=[-3,2[$ et $J=]0,5]$, donner $I\cap J$, $I\cup J$ et $\R\setminus I$.

\section*{Exercice 3.3 -- Distance}
Résoudre $|x+2|\leq3$. Écrire $x\in[4,10]$ à l'aide d'une valeur absolue.

\section*{Exercice 3.4 -- Encadrement}
Encadrer $\sqrt3$ entre deux centièmes consécutifs. Donner son arrondi au centième.

\section*{Exercice 3.5 -- Irrationalité}
En utilisant l'irrationalité de $\sqrt2$, prouver celle de $1+\sqrt2$.

\section*{Exercice 3.6 -- Chiffres significatifs}
Arrondir $0{,}004567$ à trois chiffres significatifs et $98765$ à deux chiffres significatifs. Donner une borne d'erreur absolue dans chaque cas.

\chapter{Puissances et racines carrées}
\section*{Exercice 4.1 -- Puissances}
Simplifier $2^4\,2^{-7}/2^{-1}$ et $(3^{-2})^3\,3^5$.

\section*{Exercice 4.2 -- Notation scientifique}
Calculer $(6\times10^{-3})(4\times10^7)/(3\times10^2)$ en notation scientifique.

\section*{Exercice 4.3 -- Racines}
Simplifier $\sqrt{50}+2\sqrt8-\sqrt{18}$ et $\sqrt{(-7)^2}$.

\section*{Exercice 4.4 -- Comparer}
Comparer exactement $2\sqrt7$ et $5$. Comparer $\sqrt{2+7}$ et $\sqrt2+\sqrt7$.

\section*{Exercice 4.5 -- Seuil}
Trouver le premier entier $n\geq0$ tel que $1{,}5^n>10$. Justifier la minimalité.

\section*{Exercice 4.6 -- Une preuve}
Démontrer $\sqrt{a/b}=\sqrt a/\sqrt b$ pour $a\geq0$, $b>0$, puis expliquer pourquoi on ne peut pas appliquer cette écriture à $a=b=-1$.

\chapter{Expressions, identités et formules}
\section*{Exercice 5.1 -- Développer}
Développer $(2x-3)^2$ et $(x+4)(x-4)$.

\section*{Exercice 5.2 -- Factoriser}
Factoriser $6x^2-15x$ et $(x+1)(2x-3)+4(x+1)$.

\section*{Exercice 5.3 -- Domaine conservé}
Simplifier $(x^2-9)/(x+3)$ en précisant son domaine. Peut-on attribuer une valeur à l'expression initiale en $-3$ ?

\section*{Exercice 5.4 -- Isoler}
Dans $V=\pi r^2h$, exprimer $h$ si $r>0$, puis $r$ si $h>0$, $V\geq0$ et $r\geq0$.

\section*{Exercice 5.5 -- Choisir la forme}
Écrire $x^2-8x+20$ sous la forme d'un carré plus une constante. En déduire son minimum sur $\R$ et résoudre l'égalité à $8$.

\section*{Exercice 5.6 -- Moyennes}
Approfondissement : pour $a,b>0$ de somme $12$, montrer $ab\leq36$ et trouver le cas d'égalité.

\chapter{Équations, inéquations et signes}
\section*{Exercice 6.1 -- Premier degré}
Résoudre $5x-7=2x+8$, puis $3(x-1)=3x+2$.

\section*{Exercice 6.2 -- Produit nul}
Résoudre $(2x-1)(x+4)=0$ et $x^2=12$.

\section*{Exercice 6.3 -- Quotient}
Résoudre $(x+3)/(x-1)=2$, puis $(x^2-1)/(x-1)=2$.

\section*{Exercice 6.4 -- Signes}
Résoudre $(x-3)(x+2)>0$ et $(x-3)/(x+2)\leq0$.

\section*{Exercice 6.5 -- Budget}
Un forfait coûte $14$ euros puis $3$ euros par séance. Avec un budget de $50$ euros, combien de séances au maximum peut-on payer ?

\section*{Exercice 6.6 -- Équivalence ?}
Un élève résout $x(x-2)=0$ en divisant par $x$ et conclut $x=2$. Expliquer et corriger son erreur.

\chapter{Repérage, distances et lieux géométriques}
\section*{Exercice 7.1 -- Distance et milieu}
Dans un repère orthonormé, $A(-1,2)$ et $B(5,10)$. Calculer $AB$ et les coordonnées du milieu $I$.

\section*{Exercice 7.2 -- Triangle rectangle}
Pour $A(0,0)$, $B(4,0)$ et $C(0,3)$, calculer les longueurs et justifier que le triangle est rectangle.

\section*{Exercice 7.3 -- Cercle}
Le cercle a pour équation $(x-2)^2+(y+1)^2=25$. Donner centre et rayon, puis déterminer si $(5,3)$ est sur ce cercle.

\section*{Exercice 7.4 -- Médiatrice}
Déterminer le lieu des points équidistants de $A(-2,0)$ et $B(4,0)$.

\section*{Exercice 7.5 -- Distance à un axe}
Pour $M(3,-4)$, donner les distances aux deux axes. Justifier que la distance à l'axe horizontal est minimale au point $(3,0)$.

\section*{Exercice 7.6 -- Rectangle général}
Prouver par coordonnées que les diagonales d'un rectangle ont même milieu et même longueur. Garder ses côtés $a,b>0$ comme paramètres.

\chapter{Vecteurs, translations et combinaisons}
\section*{Exercice 8.1 -- Déplacements}
$A(2,-1)$, $B(5,3)$, $C(-4,2)$. Calculer $\overrightarrow{AB}$ puis le point $D$ tel que $\overrightarrow{CD}=\overrightarrow{AB}$.

\section*{Exercice 8.2 -- Combinaison}
Pour $\vec u=(2,-3)$ et $\vec v=(-1,4)$, calculer $3\vec u-2\vec v$ et sa norme dans une base orthonormée.

\section*{Exercice 8.3 -- Chasles}
Simplifier $\overrightarrow{AB}+\overrightarrow{BC}-\overrightarrow{DC}$.

\section*{Exercice 8.4 -- Parallélogramme}
$A(1,1)$, $B(4,2)$, $C(5,6)$ sont trois sommets successifs d'un parallélogramme $ABCD$. Trouver $D$.

\section*{Exercice 8.5 -- Milieu vectoriel}
Montrer que $\overrightarrow{MA}+\overrightarrow{MB}=2\overrightarrow{MI}$ lorsque $I$ est le milieu de $[AB]$.

\section*{Exercice 8.6 -- Changement d'échelle}
Pour un vecteur non nul $\vec u$ de norme $5$, décrire $-\tfrac32\vec u$ : direction, sens et norme. Prouver la formule de norme par coordonnées.

\chapter{Colinéarité, déterminant et bases}
\section*{Exercice 9.1 -- Déterminants}
Calculer les déterminants de $(2,5)$ et $(6,15)$, puis de $(2,5)$ et $(3,7)$. Conclure.

\section*{Exercice 9.2 -- Alignement}
$A(1,-1)$, $B(3,2)$, $C(5,5)$. Démontrer l'alignement.

\section*{Exercice 9.3 -- Paramètre}
Déterminer $t$ pour que $(t,6)$ et $(2,3)$ soient colinéaires.

\section*{Exercice 9.4 -- Coordonnée nulle}
Expliquer pourquoi $(0,4)$ et $(0,-7)$ sont colinéaires sans écrire de quotient de dénominateur nul.

\section*{Exercice 9.5 -- Dans une autre base}
Exprimer $\vec w=(7,1)$ dans la base $\vec u=(1,2)$, $\vec v=(2,-1)$.

\section*{Exercice 9.6 -- Parallèles ou confondues}
$A(0,0)$, $B(2,1)$, $C(0,3)$, $D(4,5)$. Comparer les droites $(AB)$ et $(CD)$.

\chapter{Droites et systèmes linéaires}
\section*{Exercice 10.1 -- Deux points}
Trouver une équation de la droite passant par $A(1,3)$ et $B(4,9)$.

\section*{Exercice 10.2 -- Point et directeur}
Donner une équation de la droite passant par $(2,-1)$ et dirigée par $(3,2)$.

\section*{Exercice 10.3 -- Verticale}
Décrire la droite passant par $(4,-2)$ et $(4,7)$. Pourquoi la formule de pente ne s'applique-t-elle pas ?

\section*{Exercice 10.4 -- Intersection}
Trouver l'intersection de $y=3x-2$ et $y=-2x+8$.

\section*{Exercice 10.5 -- Trois systèmes}
Décrire les solutions de : $x+y=3$, $2x+2y=6$ ; puis $x+y=3$, $2x+2y=7$ ; puis $x+y=3$, $x-y=1$.

\section*{Exercice 10.6 -- Région admissible}
Approfondissement : décrire la région $x\geq0$, $y\geq0$, $2x+y\leq6$, et donner ses sommets.

\chapter{Configurations et preuves géométriques}
\section*{Exercice 11.1 -- Milieux}
Dans un triangle $ABC$, $M$ et $N$ sont les milieux de $[AB]$ et $[AC]$. Démontrer $\overrightarrow{MN}=\tfrac12\overrightarrow{BC}$.

\section*{Exercice 11.2 -- Barycentre}
Approfondissement : trouver le barycentre de $A(0,2)$ de poids $1$ et $B(6,-1)$ de poids $2$.

\section*{Exercice 11.3 -- Médianes}
Approfondissement : pour $A(0,0)$, $B(6,0)$, $C(0,9)$, trouver l'intersection des médianes et sa position sur la médiane issue de $A$.

\section*{Exercice 11.4 -- Cercle circonscrit}
Approfondissement : pour le triangle de l'exercice 7.2, trouver le centre et le rayon du cercle circonscrit.

\section*{Exercice 11.5 -- Aire}
Approfondissement : deux côtés d'un triangle mesurent $5$ et $8$, et leur angle compris est $30\degre$. Calculer l'aire et expliquer la formule utilisée.

\section*{Exercice 11.6 -- Lieu}
Approfondissement : déterminer le lieu des points équidistants de $F(0,2)$ et de l'axe des abscisses. Vérifier le point $(2,2)$.

\chapter{Fonctions : définir, représenter, modéliser}
\section*{Exercice 12.1 -- Images et antécédents}
Pour $f(x)=(x-1)^2$ sur $\R$, calculer $f(-2)$ et donner les antécédents de $4$, de $0$, puis de $-1$.

\section*{Exercice 12.2 -- Domaine}
Déterminer le domaine de $f(x)=\sqrt{3-x}/(x+2)$.

\section*{Exercice 12.3 -- Appartenance}
Le point $A(2,5)$ appartient-il à la courbe $y=x^2+1$ ? Et $B(-2,-3)$ ?

\section*{Exercice 12.4 -- Tableau insuffisant}
Deux fonctions prennent les valeurs $0,0,0$ en $-1,0,1$. Cela prouve-t-il qu'elles sont identiques ? Donner un contre-exemple.

\section*{Exercice 12.5 -- Modèle géométrique}
Une clôture de $30$ mètres entoure un rectangle. Si un côté vaut $x$, exprimer l'aire et donner le domaine des rectangles non dégénérés.

\section*{Exercice 12.6 -- Modèle discret}
Un envoi de $n$ colis coûte $u(n)=8+2{,}5n$, pour $n\in\N$. Calculer $u(4)$ et expliquer si le prix $19$ euros est réalisable.

\chapter{Fonctions affines et comparaisons}
\section*{Exercice 13.1 -- Reconstruire}
Une fonction affine vérifie $f(1)=5$ et $f(4)=14$. Déterminer son expression.

\section*{Exercice 13.2 -- Signe et variation}
Étudier le signe et les variations de $f(x)=-4x+12$.

\section*{Exercice 13.3 -- Comparer des tarifs}
Pour $x\geq0$, comparer $A(x)=30+2x$ et $B(x)=10+6x$.

\section*{Exercice 13.4 -- Une preuve}
Prouver qu'une fonction affine de pente positive est strictement croissante.

\section*{Exercice 13.5 -- Interpolation}
On suppose une évolution affine de $18$ à $30$ entre les instants $t=2$ et $t=8$. Estimer la valeur à $t=5$. Préciser l'hypothèse.

\section*{Exercice 13.6 -- Modèle physique}
Un réservoir de $75$ litres perd $2{,}5$ litres par minute. Donner la fonction de volume, son domaine utile et l'instant où il reste $20$ litres.

\chapter{Valeur absolue et fonction carré}
\section*{Exercice 14.1 -- Distances}
Résoudre $|x-1|=4$, puis $|x-1|\leq4$.

\section*{Exercice 14.2 -- Carrés}
Résoudre $x^2<16$ et $x^2\geq5$.

\section*{Exercice 14.3 -- Comparer}
Comparer sans calcul décimal $0{,}7$ et $0{,}7^2$, puis $1{,}3$ et $1{,}3^2$. Justifier par un signe.

\section*{Exercice 14.4 -- Parabole déplacée}
Étudier le minimum et les variations de $f(x)=(x+2)^2-3$ sur $\R$.

\section*{Exercice 14.5 -- Solutions et domaine}
Résoudre $(x+2)^2-3=6$, puis ne garder que les solutions dans $[0;+\infty[$.

\section*{Exercice 14.6 -- Un contre-exemple}
Réfuter « si $a<b$, alors $a^2<b^2$ ». Donner une hypothèse supplémentaire qui rende l'affirmation vraie et la démontrer.

\chapter{Inverse, racine carrée et cube}
\section*{Exercice 15.1 -- Inverse}
Résoudre $1/x=4$ et $1/x\leq-2$.

\section*{Exercice 15.2 -- Racine}
Résoudre $\sqrt{x-1}=4$, puis $\sqrt{x-1}<4$.

\section*{Exercice 15.3 -- Cube}
Résoudre $x^3=-27$ et $x^3<8$.

\section*{Exercice 15.4 -- Comparer des images}
Comparer $1/(-4)$ et $1/(-2)$, puis $\sqrt{11}$ et $\sqrt{12}$.

\section*{Exercice 15.5 -- Solution parasite}
Résoudre $\sqrt{x+6}=x$ en justifiant les conditions avant de mettre au carré.

\section*{Exercice 15.6 -- Symétrie}
Approfondissement : prouver que la courbe de la racine est l'image de la courbe du carré sur $[0,+\infty[$ par la symétrie par rapport à $y=x$.

\chapter{Variations, extrémums et optimisation}
\section*{Exercice 16.1 -- Lire des variations}
Une fonction décroît de $8$ à $-1$ sur $[-2;1]$, puis croît de $-1$ à $4$ sur $[1;5]$. Donner ses extrémums sur $[-2;5]$. Peut-on comparer $f(0)$ et $f(4)$ ?

\section*{Exercice 16.2 -- Borne non atteinte}
La fonction $f(x)=x^2$ possède-t-elle un maximum et un minimum sur $]0;2[$ ?

\section*{Exercice 16.3 -- Optimiser}
Un rectangle a pour périmètre $24$ mètres. Déterminer son aire maximale en justifiant qu'elle est atteinte.

\section*{Exercice 16.4 -- Dichotomie}
Combien d'étapes de dichotomie depuis un intervalle d'amplitude $1$ garantissent une amplitude au plus $0{,}01$ ?

\section*{Exercice 16.5 -- Grille et preuve}
Un programme cherche le maximum de $f(x)=1-(x-0{,}3)^2$ sur la grille $\{0,0{,}5,1\}$. Quel résultat trouve-t-il ? Quel est le véritable maximum sur $[0;1]$ ?

\section*{Exercice 16.6 -- Longueur approchée}
Pour $y=x^2$ sur $[0;1]$, calculer la longueur polygonale avec un seul segment, puis avec les abscisses $0,1/2,1$.

\chapter{Proportions et évolutions}
\section*{Exercice 17.1 -- Une proportion}
$54$ personnes représentent $18\,\%$ d'un groupe. Quel est son effectif ?

\section*{Exercice 17.2 -- Proportion de proportion}
$40\,\%$ des élèves prennent un bus ; parmi eux, $30\,\%$ prennent la ligne A. Quelle proportion de tous les élèves prend cette ligne ?

\section*{Exercice 17.3 -- Successions}
Une quantité augmente de $10\,\%$, puis baisse de $15\,\%$. Donner le taux global.

\section*{Exercice 17.4 -- Réciproque}
Quel taux compense exactement une baisse de $30\,\%$ ?

\section*{Exercice 17.5 -- Points ou pourcentage}
Une proportion passe de $12\,\%$ à $15\,\%$. Donner la hausse en points et la hausse relative.

\section*{Exercice 17.6 -- Pondérer}
Un groupe de $40$ personnes compte $25\,\%$ de réponses oui ; un autre de $60$ en compte $50\,\%$. Calculer la proportion globale.

\chapter{Statistiques : position et lecture critique}
\section*{Exercice 18.1 -- Indicateurs}
Pour $1,2,2,3,5,5,7,8,9,10$, calculer moyenne, médiane, quartiles, étendue et écart interquartile avec les conventions du cours.

\section*{Exercice 18.2 -- Moyenne de groupes}
Un groupe de $12$ élèves a une moyenne de $14$, un groupe de $18$ a une moyenne de $11$. Donner la moyenne globale.

\section*{Exercice 18.3 -- Transformation}
Chaque mesure d'une série de moyenne $7$ est transformée par $y=3x-2$. Quelle est la nouvelle moyenne ? Justifier.

\section*{Exercice 18.4 -- Ajout extrême}
Ajouter $100$ à la série de l'exercice 18.1. Recalculer moyenne et médiane, puis commenter.

\section*{Exercice 18.5 -- Lire un diagramme}
Une catégorie compte $45$ personnes sur $180$. Quel angle occupe-t-elle dans un diagramme circulaire, puis semi-circulaire ?

\section*{Exercice 18.6 -- Mêmes centres ?}
Comparer les séries $A=(4,5,6)$ et $B=(0,5,10)$. Peut-on les dire identiques parce qu'elles ont même moyenne et même médiane ?

\chapter{Dispersion et séries regroupées en classes}
\section*{Exercice 19.1 -- Calculer une dispersion}
Pour la série $2,4,4,6$, calculer moyenne, variance et écart-type. Contrôler la variance de deux façons.

\section*{Exercice 19.2 -- Changer d'unité}
Une série de températures a une moyenne de $18\degre$C et un écart-type de $3\degre$C. La conversion est $F=1{,}8C+32$. Calculer les deux indicateurs en degrés Fahrenheit et justifier l'absence de $32$ dans le second.

\section*{Exercice 19.3 -- Classes et histogramme}
Les classes $[0;10[$, $[10;20[$, $[20;30]$ ont pour effectifs $4,8,8$. Tracer un histogramme d'amplitudes égales et le polygone des fréquences cumulées. Estimer moyenne et médiane sous l'hypothèse de répartition uniforme dans chaque classe.

\section*{Exercice 19.4 -- Ce que le regroupement perd}
Avec les mêmes classes et effectifs, les moyennes exactes de classe sont $4$, $14$ et $26$. Déterminer la moyenne exacte globale. Peut-on en déduire l'écart-type exact ? Construire une explication.

\section*{Exercice 19.5 -- Choisir les indicateurs}
Comparer les séries $A=(10,11,12,13,14)$ et $B=(10,11,12,13,100)$. Calculer médiane, quartiles, moyenne et étendue. Expliquer pourquoi un seul indicateur serait trompeur.

\section*{Exercice 19.6 -- Décomposer une variance (facultatif)}
Deux groupes de même effectif ont des moyennes $2$ et $8$ et des variances $1$ et $4$. Calculer la moyenne et la variance de la réunion. Pourquoi la moyenne des variances ne convient-elle pas ?

\chapter{Tableaux croisés et fréquences conditionnelles}
\section*{Exercice 20.1 -- Compléter des effectifs}
Sur $200$ personnes, $120$ habitent le site A. Parmi elles, $90$ portent un casque. Sur le site B, $40$ portent un casque. Construire le tableau croisé, marges comprises.

\section*{Exercice 20.2 -- Choisir le dénominateur}
Avec ce tableau, calculer la proportion de porteurs de casque, la proportion de porteurs parmi les habitants de A et la proportion d'habitants de A parmi les porteurs. Interpréter chaque résultat.

\section*{Exercice 20.3 -- Nominal ou ordinal}
Classer les caractères suivants : canton de résidence ; satisfaction (faible, moyenne, forte) ; nombre d'enfants ; taille en centimètres. Pour lesquels une moyenne des modalités aurait-elle un sens ?

\section*{Exercice 20.4 -- Lire des données réelles}
Au 1\textsuperscript{er} janvier 2025, le canton d'Esch-sur-Alzette compte $97\,721$ femmes et $97\,370$ hommes. Le Luxembourg entier compte $339\,161$ femmes et $342\,812$ hommes (STATEC). Calculer la part des femmes dans ce canton et la part des résidentes du pays qui habitent ce canton.

\section*{Exercice 20.5 -- Filtrer avec ET, OU, NON}
Dans le tableau du casque, combien de personnes sont sur A ET portent un casque ? Sur A OU portent un casque ? Ni sur A ni porteuses de casque ? Expliquer pourquoi une addition directe échoue dans le deuxième cas.

\section*{Exercice 20.6 -- Une inversion trompeuse}
Une classe de $30$ élèves compte $18$ filles ; $12$ élèves font du théâtre, dont $9$ filles. « Les trois quarts des filles font du théâtre » est-il vrai ? Écrire l'énoncé réellement associé à $9/12$.

\chapter{Probabilités finies et dénombrement}
\section*{Exercice 21.1 -- Construire une loi}
Une variable prend les valeurs $1,2,3,4$ avec des probabilités proportionnelles à ces nombres. Déterminer la loi et la probabilité d'un résultat pair.

\section*{Exercice 21.2 -- Deux dés}
On lance deux dés équilibrés indépendants. Calculer la probabilité que la somme soit $7$, puis qu'au moins un dé donne $6$. Justifier l'univers utilisé.

\section*{Exercice 21.3 -- Réunion et intersection}
On sait $P(A)=0{,}4$, $P(B)=0{,}5$ et $P(A\cap B)=0{,}2$. Calculer $P(A\cup B)$ et la probabilité de n'avoir ni A ni B. Démontrer la formule employée en séparant trois régions disjointes.

\section*{Exercice 21.4 -- Remettre ou ne pas remettre}
Une urne contient trois boules rouges et deux bleues, indiscernables au toucher. On en tire deux. Calculer la probabilité de deux rouges avec remise, puis sans remise.

\section*{Exercice 21.5 -- Un modèle est une hypothèse}
Une pièce a donné $63$ faces en $100$ lancers. Peut-on conclure que sa probabilité de face vaut exactement $0{,}63$ ? Qu'apporterait une simulation de pièces équilibrées ?

\section*{Exercice 21.6 -- Compter un complément}
Un code comporte trois chiffres, chacun choisi indépendamment et uniformément parmi $0,\ldots,9$. Les répétitions et le zéro initial sont autorisés. Calculer la probabilité d'au moins une répétition.

\chapter{Probabilités conditionnelles et arbres}
\section*{Exercice 22.1 -- Du tableau à la probabilité}
On choisit uniformément une personne dans le tableau du casque du chapitre 20 : A avec casque $90$, A sans casque $30$, B avec casque $40$, B sans casque $40$. Calculer $P(A)$, $P_A(C)$, $P_C(A)$ et $P(A\cap C)$.

\section*{Exercice 22.2 -- Deux chemins}
Un objet provient de l'atelier A avec probabilité $0{,}6$, sinon de B. La proportion de défauts est $2\,\%$ dans A et $5\,\%$ dans B. Construire un arbre. Calculer la probabilité d'un défaut et la part des objets de A parmi les objets défectueux.

\section*{Exercice 22.3 -- Un diagnostic fictif}
Dans un modèle fictif, $1\,\%$ d'une population est malade. Le test est positif chez $90\,\%$ des malades et chez $5\,\%$ des non-malades. Construire un tableau de proportions ramenées à $10\,000$ personnes. Un positif a-t-il $90\,\%$ de probabilité d'être malade ?

\section*{Exercice 22.4 -- Sensibilité et spécificité}
Dans le modèle précédent, donner sensibilité, spécificité, taux de faux négatifs parmi les malades et taux de faux positifs parmi les non-malades. Donner aussi la probabilité d'un faux positif dans toute la population.

\section*{Exercice 22.5 -- Indépendance (facultatif)}
Pour deux événements de probabilités non nulles, on a $P(A)=0{,}4$, $P(B)=0{,}5$ et $P(A\cap B)=0{,}2$. Sont-ils indépendants ? Sont-ils incompatibles ? Justifier par les définitions.

\section*{Exercice 22.6 -- Pourquoi la multiplication ?}
On suppose $P(A)>0$. À partir de la définition d'une probabilité conditionnelle, démontrer la règle du produit le long d'un chemin. Expliquer ce qui change si $P(A)=0$.

\chapter{Algorithmique et Python}
\section*{Exercice 23.1 -- Tracer des variables}
Exécuter à la main : \begin{lstlisting}[language=Python]
s = 0
for k in range(1, 5):
    s = s + k * k
\end{lstlisting}
Donner les valeurs successives de \texttt{s}. Que changerait \texttt{range(1, 4)} ?

\section*{Exercice 23.2 -- Dernier multiple}
Écrire une fonction qui renvoie le dernier multiple d'un entier $a>0$ inférieur ou égal à un entier $b\ge0$. Proposer un test de frontière et expliquer la correction.

\section*{Exercice 23.3 -- Un encadrement certifié}
Pour encadrer $\sqrt2$ au millième, un programme cherche le plus grand entier $k$ tel que $k^2\le2\times1000^2$. Donner $k$ et l'encadrement. Pourquoi ces comparaisons entières sont-elles utiles ?

\section*{Exercice 23.4 -- Lire une fonction statistique}
Expliquer le rôle de chaque ligne, les conditions d'entrée et la différence entre le résultat renvoyé et une variance :\begin{lstlisting}[language=Python]
from math import sqrt
def dispersion(valeurs):
    n = len(valeurs)
    m = sum(valeurs) / n
    return sqrt(sum((x - m)**2 for x in valeurs) / n)
\end{lstlisting}

\section*{Exercice 23.5 -- Simuler sans conclure trop vite}
Écrire une fonction simulant la fréquence de face pour $n>0$ lancers d'une pièce équilibrée. Pourquoi deux appels peuvent-ils donner des résultats différents ?

\section*{Exercice 23.6 -- Longueur polygonale}
Écrire le principe d'un algorithme qui approche la longueur de la courbe $y=x^2$ sur $[0;1]$ avec $n$ segments. Quel résultat exact doit-il produire pour $n=1$ ? Que ne garantit pas à lui seul un tableau de valeurs croissantes ?

\chapter{Approfondir les preuves et préparer la suite}
\section*{Exercice 24.1 -- L'ordre des quantificateurs}
\textbf{Facultatif.} Comparer les affirmations : « pour tout réel $x$, il existe un réel $y$ tel que $y>x$ » et « il existe un réel $y$ tel que, pour tout réel $x$, $y>x$ ». Écrire leurs négations en français.

\section*{Exercice 24.2 -- Une somme télescopique}
\textbf{Facultatif.} Démontrer $1+3+\cdots+(2n-1)=n^2$ en utilisant $k^2-(k-1)^2$. Vérifier pour $n=4$ puis expliquer pourquoi cette vérification ne suffit pas.

\section*{Exercice 24.3 -- Une récurrence guidée}
\textbf{Facultatif.} Pour $u\ge0$ et $n\in\N$, démontrer par récurrence $(1+u)^n\ge1+nu$. En déduire que pour tout seuil $M>0$ et $u>0$, les puissances $(1+u)^n$ finissent par dépasser $M$.

\section*{Exercice 24.4 -- Une borne de fluctuation}
\textbf{Facultatif.} Pour une fréquence $F_n$ de succès lors d'épreuves indépendantes de probabilité $p$, le cours démontre $P(|F_n-p|\ge\varepsilon)\le1/(4n\varepsilon^2)$. Calculer cette borne pour $n=10\,000$ et $\varepsilon=0{,}05$. Peut-on en déduire que tout échantillon est proche de $p$ ?

\section*{Exercice 24.5 -- Trouver le trou d'une preuve}
\textbf{Facultatif.} Un texte dit : « $a=b$, donc $a^2=ab$, donc $(a-b)(a+b)=b(a-b)$, donc $a+b=b$, donc $2=1$ en prenant $a=b=1$ ». Localiser et expliquer l'erreur.

\section*{Exercice 24.6 -- Démontrer seul}
\textbf{Facultatif.} Soient $a,b>0$. Démontrer $\frac{a+b}{2}\ge\sqrt{ab}$, préciser le cas d'égalité, puis en déduire $\frac1a+\frac1b\ge\frac4{a+b}$.

\chapter{Six problèmes de synthèse}
\section*{Problème 1 -- Un jardin et son coût}
Un jardin rectangulaire s'appuie sur un mur. On dispose de $40\unite m$ de grillage pour les trois autres côtés. On note $x$ la largeur perpendiculaire au mur.
\begin{enumerate}
\item Exprimer la longueur et donner le domaine géométrique de $x$.
\item Montrer que l'aire vaut $A(x)=200-2(x-10)^2$. Déterminer son maximum.
\item Déterminer les largeurs donnant une aire d'au moins $150\unite{m^2}$.
\item Deux entreprises proposent $120+15L$ et $200+12L$ euros pour poser $L$ mètres de grillage. Comparer les devis pour $L=40$, puis pour un $L\ge0$ quelconque.
\end{enumerate}

\section*{Problème 2 -- Un triangle sous plusieurs regards}
Dans un repère orthonormé, $A(0;0)$, $B(6;0)$ et $C(2;4)$. On note I le milieu de BC et G le point de coordonnées $(8/3;4/3)$.
\begin{enumerate}
\item Calculer les longueurs des côtés et l'aire.
\item Trouver une équation de la droite AI et montrer que G y appartient.
\item Vérifier que G appartient aussi à la médiane issue de B ; préciser sa position sur chaque médiane.
\item Trouver le centre O du cercle circonscrit et son rayon.
\item Déterminer le point d'intersection des hauteurs issues de C et de A.
\end{enumerate}
Les questions 3 à 5 prolongent facultativement la partie obligatoire.

\section*{Problème 3 -- Un tarif et une expérience}
Un appareil coûte $C(n)=80+0{,}25n$ euros pour $n$ utilisations, $n$ entier strictement positif. Son coût moyen est $m(n)=C(n)/n$.
\begin{enumerate}
\item Écrire $m(n)$ à l'aide de la fonction inverse et justifier son sens de variation.
\item Déterminer le plus petit $n$ tel que $m(n)\le1$ euro.
\item Un élève trace $m(x)$ sur $]0;200]$. Expliquer la différence entre ce tracé et les valeurs autorisées par le modèle.
\item Écrire une boucle cherchant le premier entier vérifiant le seuil, et prouver qu'elle se termine.
\end{enumerate}

\section*{Problème 4 -- Comparer deux distributions}
Deux groupes de $20$ élèves ont obtenu les notes suivantes : groupe A, dix $8$ et dix $12$ ; groupe B, cinq $4$, dix $10$ et cinq $16$.
\begin{enumerate}
\item Calculer moyenne, médiane, quartiles et écart-type dans chaque groupe.
\item Représenter les diagrammes en boîte et rédiger une comparaison.
\item Le groupe A gagne uniformément $2$ points ; le groupe B gagne $10\,\%$ de sa note. Calculer les nouvelles moyennes et dispersions.
\item Peut-on affirmer que les deux transformations profitent pareillement à tous les élèves ?
\end{enumerate}

\section*{Problème 5 -- Du contrôle qualité à la décision}
Une production comporte $3\,\%$ de pièces défectueuses. Un contrôle rejette $95\,\%$ des pièces défectueuses et $4\,\%$ des pièces conformes.
\begin{enumerate}
\item Construire un arbre puis un tableau pour $100\,000$ pièces théoriques.
\item Calculer la probabilité de rejet et la proportion de pièces défectueuses parmi les rejets.
\item Calculer la proportion de pièces défectueuses parmi les pièces acceptées.
\item Un second contrôle identique permet-il de multiplier automatiquement les deux probabilités de rejet ? Expliquer l'hypothèse manquante.
\end{enumerate}

\section*{Problème 6 -- Un atelier de preuve et de données}
On étudie les populations des douze cantons dans le fichier CSV fourni.
\begin{enumerate}
\item Vérifier les totaux femmes, hommes et ensemble.
\item Démontrer que la proportion nationale de femmes est une moyenne des proportions cantonales pondérée par les populations des cantons.
\item Expliquer pourquoi la moyenne simple des douze proportions ne donne pas forcément cette proportion nationale.
\item Programmer la construction du tableau Esch / autres cantons croisé avec femmes / hommes, puis vérifier ses marges.
\item \textbf{Facultatif.} Démontrer pour deux nombres $u,v$ et $0\le t\le1$ que $tu+(1-t)v$ est compris entre $u$ et $v$. Relier cela à la question 2.
\end{enumerate}

\end{document}
